\subsection*{电势的叠加、梯度与等势面}

静电力是保守力，可定义电势差
\[
 \varphi_B-\varphi_A=-\int_A^B\vec E\cdot\dd\vec l.
\]
电势是单位正电荷的势能，只由场源和位置决定；势能 $U=q\varphi$ 还与被放入的电荷有关。通常
把无穷远取为零势，但无限长线电荷或无限大平面的势在无穷远发散，此时只能选有限参考点讨论
电势差。

点电荷电势为 $q/(4\pi\varepsilon_0r)$，多个点电荷直接作代数和。连续分布为
\[
 \varphi(\vec r)=\frac1{4\pi\varepsilon_0}
 \int\frac{\dd q'}{|\vec r-\vec r'|}.
\]
先求标量电势再取梯度，常比直接积分矢量电场容易：
\[
 \vec E=-\nabla\varphi.
\]
负号表示电场指向电势下降最快的方向。等势面与电场线正交，沿等势面移动电荷时电场不做功；
等势面越密，法向电势变化越快，场强越大。

\subsection*{电荷系统能量与微观能标}

把点电荷逐个从无穷远缓慢搬到指定位置，外力做功储存在电荷系统中：
\[
 U=\frac12\sum_{i\ne j}
 \frac1{4\pi\varepsilon_0}\frac{q_iq_j}{r_{ij}}
 =\frac12\sum_iq_i\varphi_{\rm others}(\vec r_i).
\]
因子 $1/2$ 用来避免每一对相互作用重复计算。连续分布相应写成
$U=\tfrac12\int\rho\varphi\,\dd V$；点电荷自身的发散自能不包含在普通的粒子间相互作用能中。
求系统所受力时，可在保持其他独立坐标不变的条件下使用 $F_x=-\partial U/\partial x$。

电子伏特是能量单位：带一个元电荷的粒子通过 $1\,\mathrm V$ 电势差，能量改变
$1\,\mathrm{eV}=1.602\times10^{-19}\,\mathrm J$。原子、核与高能过程跨越的尺度很大，用 eV、keV、
MeV 表示比焦耳直观。电势高低本身不能直接判断带负电粒子的运动方向，应始终从
$\Delta K=-\Delta U=-q\Delta\varphi$ 判断。

氢原子的经典圆轨道模型把库仑力作为向心力，再配合 Bohr 角动量量子化，可得到能级按
$-1/n^2$ 分布；Rutherford 散射则用正电原子核对 $\alpha$ 粒子的库仑排斥解释大角偏转。经典轨道
图像不能替代量子理论，但它清楚展示了电势能、动能和特征尺度如何由基本常量组合出来。

在远离局域电荷分布处，电势也可作多极展开。总电荷不为零时最低阶为 $Q/r$；总电荷为零但偶极矩
不为零时，最低阶为 $\vec p\cdot\hat r/r^2$；若偶极矩也为零，四极项才主导。每高一阶都多一个
$1/r$ 的衰减，因此观察距离越远，越只看得到最低非零多极矩。

